\section*{第 \ref{ch:b3:mantle-convection} 章——地幔对流与板块驱动力}

\begin{solution}{exo:b3:mantle-convection:1}
$\mathrm{Ra} = 3500\times9.9\times\num{3e-5}\times1000\times(\num{6.7e5})^3/
(10^{-6}\times\num{5e20}) = \num{6.25e5}$，约为临界值
657.5 的 950 倍：该层发生强烈对流。
\end{solution}

\begin{solution}{exo:b3:mantle-convection:2}
$\tau_M = 10^{21}/\num{7.7e10} = \qty{1.30e10}{s}$，约 \qty{412}{\text{年}}。
潮汐持续数小时，遇到弹性地幔；冰盖熔化，
历时数千年，板片下沉历时数百万年，
均遇到黏性地幔。
\end{solution}

\begin{solution}{exo:b3:mantle-convection:3}
$\tau = 4\pi\times10^{21}/(3300\times9.8\lambda)$：\qty{3.89e11}{s} =
对 \qty{1000}{km}，为 \qty{12.3}{ka}；对 \qty{3000}{km}，为 \qty{4.1}{ka}。
荷载越宽，回弹越快。
\end{solution}

\begin{solution}{exo:b3:mantle-convection:4}
$g\rho_m\alpha(T_m - T_0)\kappa t = 9.8\times3300\times\num{3.1e-5}\times1450
\times\num{0.8e-6}\times50\times\num{3.156e13} = \qty{1.835e12}{N.m^{-1}}$；括号项为
$1 + 2\times3300\times\num{3.1e-5}\times1450/(\pi\times2270) = 1.0416$，
故 $F_{rp} = \qty{1.91e12}{N.m^{-1}}$，比板片拉力小 24 倍。
\end{solution}

\begin{solution}{exo:b3:mantle-convection:5}
(a) $v = 3300\times9.8\times\num{3e-5}\times300\times(\num{e4})^2/10^{21} =
\qty{2.91e-11}{m.s^{-1}}$；穿越 \qty{10}{km} 所需的上升时间为 $d/v =
\qty{3.44e14}{s}$，\qty{10.9}{Ma}，扩散时间为 $d^2/\kappa =
\qty{e14}{s}$，\qty{3.2}{Ma}。二者之比为 $(d^2/\kappa)/(d/v) =
vd/\kappa = \rho g\alpha\Delta Td^3/(\kappa\eta) = \mathrm{Ra} = 0.29$。(b) 先冷却的是
小块：小块不发生对流。(c) $\mathrm{Ra} \propto d^3$，故 $d =
10\times(657.5/0.291)^{1/3} = \qty{131}{km}$。
\end{solution}

\begin{solution}{exo:b3:mantle-convection:6}
$2\mathrm{Ra}_c = 1315$ 且 $\mathrm{Ra}/2\mathrm{Ra}_c = 7604$，其立方根
为 19.66。$\delta = 2891/19.66 = \qty{147}{km}$，$\mathrm{Nu} = 19.66/2 =
9.8$，$u = (10^{-6}/\num{2.891e6})\times19.66^2 = \qty{1.34e-10}{m.s^{-1}} =
\qty{0.42}{cm/\text{年}}$，且 $q = 9.8\times3.1\times1450/\num{2.891e6} =
\qty{15}{mW.m^{-2}}$。
\end{solution}

\begin{solution}{exo:b3:mantle-convection:7}
$u = 0.09/\num{3.156e7} = \qty{2.85e-9}{m.s^{-1}}$，则 $B = 2270\times\num{8e8}
\times\num{2.85e-9} = \qty{5.18e3}{kg.s^{-1}} = \qty{5.2}{Mg/s}$，约为
夏威夷的五分之三。热流为 $1170\times5178/\num{3.1e-5} =
\qty{1.95e11}{W}$，即 \qty{0.2}{TW}，或占地球
热损失的 \qty{0.42}{\percent}。
\end{solution}

\begin{solution}{exo:b3:mantle-convection:8}
$v = 2\times30\times9.8\times(\num{4e5})^2/(9\times10^{21}) =
\qty{1.05e-8}{m.s^{-1}} = \qty{0.33}{m/\text{年}}$；\qty{2000}{km} 所需
时间为 $\num{2e6}/\num{1.05e-8} = \qty{1.9e14}{s}$，约 \qty{6}{Ma}。
\end{solution}

\begin{solution}{exo:b3:mantle-convection:9}
$w = 300e^{-10/4.4} = \qty{30.9}{m}$，则 $\dot w = w/\tau = 30.9/4400 =
\qty{7.0}{mm/\text{年}}$，为乌梅奥附近实测速率的三分之二：一个均匀
半空间在该 $\tau$ 下的结果处于合理范围，但真实回弹
更快，或其凹陷更深。
\end{solution}

\begin{solution}{exo:b3:mantle-convection:10}
(a) 取 $w = 2\rho_m\alpha\Delta T\sqrt{\kappa t/\pi}/(\rho_m - \rho_w)$，
\[
  \tfrac12(\rho_m - \rho_w)gw^2 = \frac{2g\rho_m^2\alpha^2\Delta T^2\kappa t}
  {\pi(\rho_m - \rho_w)} = g\rho_m\alpha\Delta T\kappa t\times
  \frac{2\rho_m\alpha\Delta T}{\pi(\rho_m - \rho_w)} .
\] (b) 分部积分，
$\int_0^\infty x\operatorname{erfc}x\,\mathrm dx = [\tfrac{x^2}2\operatorname{erfc}
x]_0^\infty + \int_0^\infty\tfrac{x^2}2\,\tfrac2{\sqrt\pi}e^{-x^2}\,\mathrm dx =
0 + \tfrac1{\sqrt\pi}\cdot\tfrac{\sqrt\pi}4 = \tfrac14$。(c) 在 \qty{100}{Ma} 时，
$\rho_mg\alpha\int z\theta = \qty{3.67e12}{N.m^{-1}}$，并取 $w =
\qty{3705}{m}$，$\tfrac12(\rho_m - \rho_w)gw^2 = \qty{1.53e11}{N.m^{-1}}$：水的项
仅占该力的 \qty{4}{\percent}。
\end{solution}

\begin{solution}{exo:b3:mantle-convection:11}
$\ln f = 3\ln(\pi^2 + k^2) - 2\ln k$，其导数 $6k/(\pi^2 + k^2) - 2/k$
当 $3k^2 = \pi^2 + k^2$ 时为零，即 $k^2 = \pi^2/2$。此时 $f = (3\pi^2/2)^3/
(\pi^2/2) = 27\pi^4/4$。波长 $2\pi/k = 2\sqrt2\,d = 2.83\,d$ 包含
一对反向旋转的对流环流，每个宽 $\sqrt2\,d = 1.41\,d$。
\end{solution}

\begin{solution}{exo:b3:mantle-convection:12}
(a) $u = \qty{1.90e-9}{m.s^{-1}}$，$\eta u/h = 10^{20}\times\num{1.90e-9}/
\num{1.5e5} = \qty{1.27}{MPa}$，在 \qty{6000}{km} 长度上合计为 \qty{7.6e12}{N.m^{-1}}。
(b) 阻力为\omterm{prop:b3:mantle-convection:ridge-push}{洋脊推力}的两倍、板片拉力的六分之一：单独的阻力可抵抗\omterm{prop:b3:mantle-convection:ridge-push}{洋脊推力}，却不能抵抗板片拉力。
(c) 由 $\eta uL/h = \num{3.8e12} +
\num{4.6e12}$ 可得 $u = \num{8.4e12}\times\num{1.5e5}/(\num{6e6}\times10^{20})
= \qty{2.1e-9}{m.s^{-1}} = \qty{6.6}{cm/\text{年}}$，即快速大洋
板块的速度。
\end{solution}

\begin{solution}{pb:b3:mantle-convection:1}
\textbf{1.} $p = 1000\times9.8\times300 = \qty{2.94}{MPa}$。

\textbf{2.} $w = \rho_wh/\rho_m = 300\times1000/3300 = \qty{90.9}{m}$。地壳
整体下沉；从荷载下方流出又流回的
是地幔，因此该柱体由深部被排开的地幔来平衡。

\textbf{3.} $\pi\times(\num{1.25e5})^2 = \qty{4.91e10}{m^2}$，水
重 $\num{4.91e10}\times300\times1000 = \qty{1.47e16}{kg}$。

\textbf{4.} 岸线形成时为水平线，此时盆底在中心处下沉最深；此后中心比边缘抬升得更多，因此岸线现在
向中心爬升，共升高 \qty{77.2}{m}，跨度 \qty{125}{km}：
\qty{0.62}{m/km}。
\qty{0.62}{m/km}.

\textbf{5.} $\lambda/2 = \qty{250}{km}$，于是 $\lambda = \qty{500}{km}$，且 $k =
2\pi/\num{5e5} = \qty{1.257e-5}{m^{-1}}$。

\textbf{6.} $1/k = \lambda/2\pi = \qty{79.6}{km}$。

\textbf{7.} $w = 90.9 - \text{隆升}$：79.4、69.3、52.9、40.4、26.9 和
\qty{13.7}{m}。

\textbf{8.} 由 $w = w_0e^{-t/\tau}$ 得 $\ln w = \ln w_0 - t/\tau$，
这是斜率为 $-1/\tau$ 的直线。

\textbf{9.} $\tau = (14 - 1)/\ln(79.4/13.7) = 13/1.7571 = \qty{7.40}{ka}$。

\textbf{10.} $w(6) = 79.4\,e^{-5/7.40} = \qty{40.4}{m}$，隆升量为
\qty{50.5}{m}，与表中一致。

\textbf{11.} 尚余 \qty{13.7}{m}；$\dot w = 13.7/7400 = \qty{1.85}{mm/\text{年}}$。

\textbf{12.} 全部 \qty{13.7}{m}，并越来越慢地逼近。

\textbf{13.} 速率按 $e^{-t/\tau}$ 下降：$t = 7.40\ln(1.85/0.5) =
7.40\times1.309 = \qty{9.7}{ka}$，约多出 9700 年。

\textbf{14.} $\tau = 7399\times\num{3.156e7} = \qty{2.335e11}{s}$，且 $\eta =
\tau\rho_mg\lambda/(4\pi) = \num{2.335e11}\times3300\times9.8\times\num{5e5}/
12.566 = \qty{3.0e20}{Pa.s}$。

\textbf{15.} 略低于芬诺斯坎底亚的取值范围下限。波长为 \qty{500}{km} 的荷载下方的流动
在 \qty{80}{km} 的深度内衰减：
回弹所取样的是最浅、最弱的地幔，它位于
实际地球中岩石圈的下方，并部分位于其中。

\textbf{16.} $\tau = 4\pi\times\num{3.0e20}/(3300\times9.8\times\num{3e6}) =
\qty{3.9e10}{s} = \qty{1.2}{ka}$。在如此宽的荷载下方，流动可深达
\qty{477}{km} 甚至更深，朝向黏度高十倍的下地幔，
这使流动减慢。

\textbf{17.} 它以挠曲方式承担窄荷载的一部分：平衡时的下沉
较均衡值 \qty{90.9}{m} 更浅、更宽，且部分
回弹是弹性的、即时完成的。用一个均匀黏性半空间拟合
隆升数据，所得黏度便混合了两种效应。

\textbf{18.} $\tau_M = \num{3.0e20}/\num{7.7e10} = \qty{3.9e9}{s}$，约
\qty{124}{\text{年}}，比 $\tau$ 短六十倍：在数千年的时间尺度上，地幔
其响应表现为黏性流体。

\textbf{19.} $1 - e^{-50/7.40} = 0.9988$：已达到平衡下沉量的
\qty{99.9}{\percent}，假设成立。

\textbf{20.} $\eta \propto \lambda$：$\lambda$ 减半则 $\eta$ 减半，且
$\log_{10}\eta$ 下降 $\log_{10}2 = 0.30$。

\textbf{21.} $\mathrm{Ra} = 3700\times9.9\times\num{3.1e-5}\times1450\times
(\num{6.7e5})^3/(10^{-6}\times\num{3.0e20}) = \num{1.65e6}$。

\textbf{22.} $\num{1.65e6}/657.5 \approx 2500$ 倍于临界值。

\textbf{23.} $(\mathrm{Ra}/2\mathrm{Ra}_c)^{1/3} = (1254)^{1/3} = 10.78$：$\delta
= 670/10.78 = \qty{62}{km}$，且 $\mathrm{Nu} = 5.4$。

\textbf{24.} $u = (10^{-6}/\num{6.7e5})\times10.78^2 = \qty{1.73e-10}{m.s^{-1}}
= \qty{0.55}{cm/\text{年}}$，比板块慢两至十四倍。

\textbf{25.} $q = 5.4\times3.1\times1450/\num{6.7e5} = \qty{36}{mW.m^{-2}}$。该层
未计入地幔内部产生的热量，以及来自
地核与下地幔的热量；此标度律还舍去了量级为一的常数，
而板块自洋脊起以半空间方式冷却，在年轻洋底
上散失的热量远多于平均边界层的情形。

\textbf{26.} $\eta = \qty{3.0e20}{Pa.s}$：$\log_{10}\eta =
\textbf{\num[text-series-to-math=true]{20.48}}$。
\end{solution}
